

\paragraph{Simulation setup.} 
We evaluate the WAV on the LIBERO benchmark~\citep{liu2023libero}, a comprehensive suite for multi-task and lifelong robotic manipulation.
LIBERO consists of four task suites designed to probe different aspects of generalization.
\textbf{Spatial} evaluates spatial reasoning by varying scene layouts while keeping object identities fixed.
\textbf{Object} tests object-level generalization by changing object instances within a fixed environment.
\textbf{Goal} assesses goal-conditioned behavior by varying task goals under the same scene configuration.
Finally, \textbf{Long} includes long-horizon compositional tasks involving diverse objects, layouts, and goals, posing significant challenges for temporal consistency and multi-step reasoning.
A single WAV model is deployed across four suites.
\begin{table}[htbp]
\caption{Performance comparison on the LIBERO benchmark.}
\label{tab:libero_only}
\centering
\small
\resizebox{0.8\columnwidth}{!}{
\begin{tabular}{lcccccc}
\toprule
\multirow{2}{*}{\centering\textbf{Models}} 
& \multirow{2}{*}{\textbf{Params}}
& \multicolumn{4}{c}{\textbf{LIBERO}} 
& \multirow{2}{*}{\textbf{Avg.}} \\
\cmidrule(lr){3-6}
& & \textbf{Spatial} & \textbf{Object} & \textbf{Goal} & \textbf{Long} & \\
\midrule
Diffusion Policy \citet{chi2023diffusion} & - & 78.3 & 92.5 & 68.3 & 50.5 & 72.4 \\
Octo \citet{team2024octo} & - & 78.9 & 85.7 & 84.6 & 51.1 & 75.1 \\
OpenVLA \citet{kim2024openvla} & 7b & 84.9 & 88.4 & 79.2 & 53.7 & 76.5 \\
SpatialVLA \citet{qu2025spatialvla} & 4b & 88.2 & 89.9 & 78.6 & 55.5 & 78.1 \\
DiT Policy \citet{hou2024diffusion} & - & 84.2 & 96.3 & 85.4 & 63.8 & 82.4 \\
$\pi_{0}$-FAST \citet{pertsch2025fast} & 3b & 96.4 & 96.8 & 88.6 & 60.2 & 85.5 \\
GR00T-N1 \citet{bjorck2025gr00t} & 2b & 94.4 & 97.6 & 93.0 & 90.6 & 93.9 \\
$\pi_{0}$ \citet{black2024pi_0} & 3b & 96.8 & 98.8 & 95.8 & 85.2 & 94.2 \\
VLA-0 \citet{goyal2025vla} & 3b & 97.0 & 97.8 & 96.2 & 87.6 & 94.7 \\
Discrete Diffusion VLA \citet{liang2025discrete} & 7b & 97.2 & 98.6 & 97.4 & 92.0 & 96.3 \\
MemoryVLA \citet{shi2025memoryvla} & 7b & 98.4 & 98.4 & 96.4 & 93.4 & 96.5 \\
$\pi_{0.5}$ \citet{intelligence2025pi_} & 3b & 98.8 & 98.2 & 98.0 & 92.4 & 96.8 \\
OpenVLA-OFT \citet{kim2025fine} & 7b & 97.6 & 98.4 & 97.9 & 94.5 & 97.1 \\
VLA-Adapter \citet{wang2025vla} & 0.5b & 97.8 & 99.2 & 97.2 & \textbf{95.0} & 97.3 \\
\midrule
\rowcolor{gray!15}
\multicolumn{7}{l}{\textit{World Model}} \\
WorldVLA \citet{cen2025worldvla} & 7b & 85.6 & 89.0 & 82.6 & 59.0 & 79.1 \\
CoT-VLA \citet{zhao2025cot} & 7b & 87.5 & 91.6 & 87.6 & 69.0 & 81.1 \\
FlowVLA \citet{zhong2025flowvla} & 8.5b & 93.2 & 95.0 & 91.6 & 72.6 & 88.1 \\
DreamVLA \citet{zhang2025dreamvla} & 3b & 97.5 & 94.0 & 89.5 & 89.5 & 92.6 \\
UD-VLA \citet{chen2025unified} & - & 94.1 & 95.7 & 91.2 & 89.6 & 92.7 \\
RynnVLA-002-Discrete \citet{cen2025rynnvla} & 7b & 94.2 & 96.8 & 94.6 & 87.6 & 93.3 \\
UniVLA \citet{bu2025univla} & 7b & 95.4 & 98.8 & 93.5 & 94.0 & 95.5 \\
$\mathcal{F}_1$\citet{lv2025f1} & 4.5b & 98.2 & 97.8 & 95.4 & 91.3 & 95.7 \\
GE-ACT \citet{liao2025genie} & 2b & 98.2 & 97.6 & 95.8 & 94.4 & 96.5 \\
\midrule
\rowcolor{lightblue}
\textbf{WAV (Ours)} & 2.2b & \textbf{99.6} & \textbf{100.0} & \textbf{98.6} & 94.4 & \textbf{98.1} \\
\qquad W/o Latent Trajectory Planning &  & 99.0 & 99.6 & 95.0 & 91.8 & 96.4 \\
\bottomrule
\end{tabular}
}
\end{table}

\paragraph{Results.} 
Tab.~\ref{tab:libero_only} reports results on the LIBERO benchmark. 
Our method achieves state-of-the-art (or competitive) performance across all task suites, with an overall average score of 98.1. 
These results highlight its effectiveness in handling long-horizon and compositional tasks, where efficient trajectory-level inference is critical.
To assess the contribution of implicit planning, we perform an ablation by removing latent trajectory planning. 
This leads to an average performance drop of 1.7 points, with the largest degradation on the \textbf{Long} suite, decreasing from 94.4 to 91.8. 
This trend underscores the importance of trajectory-level planning for robust and generalizable task execution.
Notably, the performance gap widens with increasing task horizon, suggesting that implicit planning is particularly beneficial for mitigating compounding errors in long-horizon scenarios.

